\phantomsection\label{sec:teaching-reflections}
\section*{Reflections on Teaching}

\subsection*{Teaching Philosophy}
In my teaching, the Problem Based Learning model engages students with real-world problems and guides them to develop solutions that combine theory with practice.
I have found that students learn best when they see why a topic matters and when they take ownership of their learning, and this shapes how I design courses and supervise projects.
My goal is to help students graduate with both technical depth and the ability to keep learning on their own.

\subsection*{Course Development}
For the AI Systems \& Infrastructure course, I prepared all materials from scratch, writing over ten modules as blog-style posts.
Blog-style literature is easy to keep up to date, accessible, and tailored to the course, and I wrote it to be professional, comprehensive, intuitive, and lightly humorous.
I designed the material with the awareness that most students have little prior knowledge of the underlying systems, so it conveys the ideas in a relatable and intuitive way.

\vspace{\entryskip}
Each module pairs a concrete problem that practitioners face with hands-on exercises, and the exercises build on one another.
Students first learn HTTP protocols and build a command-line chatbot that consumes AI APIs, then reverse roles to build their own image classification API server with authentication, user management backed by a database, and rate limiting.
The next module has them package that server into Docker containers, and a later module guides them to deploy it on physical hardware such as a Raspberry Pi or an Nvidia Jetson, covering architecture compatibility and remote access.
The course ends with a mini project where groups combine these skills to deploy a versatile AI system on hardware they control.
The lectures stay inspirational and motivational, easy to follow within the limited time, and serve as an entry point for a deeper dive into each topic.
The oral exam checks whether students can solve a real-world problem with confidence, rather than recite syntax or name concepts.

\vspace{\entryskip}
The course was well received.
It was the second iteration, following a first that was poorly received, and the materials from that first iteration were unusable by my standards, so I rebuilt the course from scratch.
The official semester report was very positive, especially about the literature, with praise appearing even in the section for critical feedback.

\subsection*{Supervision Approach}
I supervise Bachelor's and Master's semester projects, and I write project proposals that are visually engaging and set a clear direction defined by a practical problem.
My supervision follows the same Problem Based Learning style.
I start students off with relatable, practical problems and real-world examples, focus on the key considerations behind a well planned project and the reasoning behind them, and give students enough resources to work out how to apply those considerations themselves.

\vspace{\entryskip}
I try to be accessible and supportive while leaving students room to make their own decisions.
I offer guidance on technical approaches and help when they hit obstacles, and I encourage them to own the direction of their project.
When I proposed the A Camera That Sees Behind topic on removing objects from images with a lightweight model deployable on embedded devices, groups extended the original proposal from their own analysis, one toward GDPR-compliant privacy and one toward the constraints of running models on embedded hardware.
Students have described my supervision as helpful and accessible, and I have watched their confidence grow as they take on deployment challenges on their own.

\subsection*{Teaching Plan}
I am satisfied with these practices and intend to keep them while revising along the way.
Going forward, I will hold to the three core principles and to the specific practices I have found effective: blog-style literature, inspirational and engaging lectures, transparency about course and exam formats and expectations, and a balanced, progressive arrangement of modules.
I am confident that I can teach any computer science course well with this approach, and I am interested in developing new courses around spatiotemporal data mining and AI infrastructure.

\vspace{\entryskip}
I am also preparing to supervise at the Master's and PhD level, carrying over the practices that matter across semester projects, research supervision, and computer science engineering in general: a clear definition of the problem, a clear path to a convincing conclusion, critical thinking about design and implementation choices, and quality in both the report and the presentation.
Through teaching grounded in research and engaged supervision, I want to help produce graduates who think critically about technology and apply their knowledge to real-world problems.
